Para evitar estos problemas sin renunciar a la IA, algunas instituciones han diseñado sistemas que siguen reglas pedagógicas muy estrictas. El referente más claro es el CS50 Duck de la Universidad de Harvard. Este asistente está integrado directamente en CS50.dev, una versión de Visual Studio Code adaptada para la nube mediante GitHub Codespaces. No funciona como un simple chat abierto, sino como un sistema complejo que utiliza una arquitectura RAG para conectar las dudas del alumno con el temario oficial del curso y filtros de seguridad que bloquean intentos de inyección de prompts. Toda su lógica operativa se gestiona desde un backend centralizado que utiliza archivos de configuración YAML para ajustar dinámicamente el comportamiento de GPT-4, asegurando que el modelo nunca entregue la solución directa, sino pistas y preguntas de seguimiento que fomenten el razonamiento independiente y la técnica del rubberducking\cite{ malan2026cs50ai}.

Para gestionar los costes de la API comercial y potenciar el aprendizaje profundo, Harvard implementó un ingenioso sistema de "corazones" que limita las consultas de los estudiantes, obligándoles a reflexionar antes de formular sus dudas para no agotar sus recursos. Aunque este enfoque es pedagógicamente brillante, su dependencia de modelos en la nube genera gastos operativos elevados y plantea retos sobre la privacidad al enviar datos a servidores externos\cite{liu2024cs50}.

En una línea parecida el asistente CodeAid ha validado un modelo de ayuda estructurada en clases masivas. A diferencia de un chat convencional este software ofrece un menú de herramientas específicas para el alumno como la explicación de conceptos técnicos o la resolución de dudas teóricas. Para la revisión del código el sistema habilita una caja de texto donde el usuario puede pegar sus líneas pero sin capacidad de ejecución directa. A partir de ahí el programa sigue un proceso interno basado en una técnica de \textit{few-shot prompting}, donde se entrena al modelo con ejemplos previos de interacciones pedagógicas para asegurar que el tono sea siempre de apoyo. Bajo esta arquitectura la IA primero genera la solución correcta de forma invisible para garantizar la precisión pero después la traduce en un pseudocódigo que orienta al estudiante sin darle nunca la respuesta final. Este método tiene el beneficio de forzar al usuario a que reflexione sobre la lógica del ejercicio y tenga que escribir él mismo cada línea en su editor lo que asegura que el proceso de aprendizaje sea efectivo y se evite el simple copiado de soluciones. Además esta plataforma también anonimiza los datos personales al igual que hace el sistema de Harvard aunque no cuenta con una protección tan estricta contra los ataques de inyección de instrucciones\cite{kazemitabaar2024codeaid}.